% Generated draft: requires mathematical and source review.
\documentclass[11pt]{article}
\usepackage[margin=25mm]{geometry}
\usepackage{fontspec}
\setmainfont{Latin Modern Roman}
\usepackage{amsmath,amssymb,mathtools,unicode-math}
\setmathfont{Latin Modern Math}
\usepackage{enumitem,needspace,xurl,hyperref}
\setlength{\parindent}{0pt}
\setlength{\parskip}{5pt}
\setlength{\emergencystretch}{3em}
\newcommand{\N}{\mathbb{N}}
\newcommand{\Z}{\mathbb{Z}}
\newcommand{\Q}{\mathbb{Q}}
\newcommand{\R}{\mathbb{R}}
\newcommand{\C}{\mathbb{C}}
\newcommand{\F}{\mathbb{F}}
\newcommand{\A}{\mathbb{A}}
\newcommand{\PP}{\mathbb{P}}
\newcommand{\E}{\mathbb{E}}
\newcommand{\eps}{\varepsilon}
\newcommand{\dd}{\,\mathrm{d}}
\begin{document}
\section*{UnsolvedMath ID: 20001143}
\textbf{Source problem code:} catalogue entry
\subsection{Definitions and mathematical statement}
Use the normalization explicitly recovered in the supplied source verification:
\[
 \mathcal C_4=\left\{\sum_{j=0}^{\infty}\frac{a_j}{4^j}:a_j\in\{0,1\}\right\}\subseteq[0,4/3].
\]
For $\alpha\in\R$ let $E_\alpha=\mathcal C_4+\alpha\mathcal C_4=\{x+\alpha y:x,y\in\mathcal C_4\}$. T. Tao asks to characterize exactly the slopes $\alpha$ for which one-dimensional Lebesgue measure $\mathcal L^1(E_\alpha)$ is zero. The workshop remarks that this is known for Lebesgue-almost every real $\alpha$.

The attached aid reports the complete classification
\[
 \mathcal L^1(E_\alpha)=0
 \iff
 \begin{cases}
 \alpha=0,\quad\text{or}\\
 \alpha\notin\Q,\quad\text{or}\\
 \alpha=p/q\ne0,\ \gcd(|p|,q)=1,\ q>0,
 \quad v_2(|p|)-v_2(q)\equiv0\pmod2.
 \end{cases}
\]
Here $v_2(a)$ is the exponent of $2$ in a positive integer $a$; the gap may be negative, with parity taken in $\Z$. Equivalently, a nonzero reduced rational has positive measure exactly when this valuation gap is odd. The aid further reports that these positive-measure sets contain a nondegenerate interval; irrational slopes have Hausdorff dimension one despite zero measure; and nonzero rational zero-measure cases have an exact self-similar overlap and Hausdorff dimension less than one.

For the normalized positive rational subcase $p=2u>0$ with $u,q$ odd and $\gcd(p,q)=1$, it also records the bound $\mathcal L^1(E_{p/q})\geq4/q$. This uses the displayed index-$0$ convention; changing the starting index rescales the measure. The aid describes reductions by sign, inversion, and multiplication of the slope by $4$, followed by complete-residue and Fourier certificates for rational slopes. No proof or full published classification reference is supplied here, so the characterization and dimensional refinements remain recorded claims requiring independent verification.
\subsection{Short English statement}
Characterize zero-measure projections of the four-corner Cantor product. The supplied classification says all irrational slopes have zero measure, while a nonzero rational has positive measure exactly when its reduced numerator/denominator 2-adic valuation gap is odd. The starting-index normalization is explicit.
\subsection{Sources}
\begin{itemize}
\item AIM list \emph{Recent trends in additive combinatorics}, problem 4.12 attributed to T. Tao, referring to the definition in problem 4.11: \url{https://aimath.org/WWN/additivecomb/additivecomb.pdf}. The supplied verification says the official TeX starts the Cantor series at index zero; this draft did not retrieve that TeX or the PDF.
\item Local registry \texttt{data/unsolvedmath/problems.json}, permanent ID 20001143, exact source code \texttt{catalogue entry}; canonical location \texttt{aim-combinatorics-notes.json notes[267]}.
\item Attached attempt 1, full solution, problem status at run solved, rigor proved, low novelty confidence. It invokes a published classification but supplies no author/title/full bibliographic reference. It is machine-generated, not peer reviewed, requiring independent verification before citation.
\end{itemize}
\end{document}
